\documentclass[10pt,a4paper]{amsart}
\usepackage{amsmath,amssymb,amsthm,graphicx}
% For commutative diagrams we use Xy-pic
\usepackage[all]{xy}

% We use 2cell for 2-commutative diagrams.
\xyoption{2cell}
% Theorem environments.
%
\theoremstyle{plain}
\newtheorem{theorem}[subsection]{Theorem}
\newtheorem{proposition}[subsection]{Proposition}
\newtheorem{lemma}[subsection]{Lemma}

\theoremstyle{definition}
\newtheorem{definition}[subsection]{Definition}
\newtheorem{example}[subsection]{Example}
\newtheorem{exercise}[subsection]{Exercise}
\newtheorem{situation}[subsection]{Situation}

\theoremstyle{remark}
\newtheorem{remark}[subsection]{Remark}
\newtheorem{remarks}[subsection]{Remarks}

\numberwithin{equation}{subsection}

% Macros
%
\def\lim{\mathop{\mathrm{lim}}\nolimits}
\def\colim{\mathop{\mathrm{colim}}\nolimits}
\def\Spec{\mathop{\mathrm{Spec}}}
\def\Hom{\mathop{\mathrm{Hom}}\nolimits}
\def\Ext{\mathop{\mathrm{Ext}}\nolimits}
\def\SheafHom{\mathop{\mathcal{H}\!\mathit{om}}\nolimits}
\def\SheafExt{\mathop{\mathcal{E}\!\mathit{xt}}\nolimits}
\def\Sch{\mathit{Sch}}
\def\Mor{\mathop{\mathrm{Mor}}\nolimits}
\def\Ob{\mathop{\mathrm{Ob}}\nolimits}
\def\Sh{\mathop{\mathit{Sh}}\nolimits}
\def\NL{\mathop{N\!L}\nolimits}
\def\CH{\mathop{\mathrm{CH}}\nolimits}
\def\proetale{{pro\text{-}\acute{e}tale}}
\def\etale{{\acute{e}tale}}
\def\QCoh{\mathit{QCoh}}
\def\Ker{\mathop{\mathrm{Ker}}}
\def\Im{\mathop{\mathrm{Im}}}
\def\Coker{\mathop{\mathrm{Coker}}}
\def\Coim{\mathop{\mathrm{Coim}}}

% Boxtimes
%
\DeclareMathSymbol{\boxtimes}{\mathbin}{AMSa}{"02}

%
% Macros for moduli stacks/spaces
%
\def\QCohstack{\mathcal{QC}\!\mathit{oh}}
\def\Cohstack{\mathcal{C}\!\mathit{oh}}
\def\Spacesstack{\mathcal{S}\!\mathit{paces}}
\def\Quotfunctor{\mathrm{Quot}}
\def\Hilbfunctor{\mathrm{Hilb}}
\def\Curvesstack{\mathcal{C}\!\mathit{urves}}
\def\Polarizedstack{\mathcal{P}\!\mathit{olarized}}
\def\Complexesstack{\mathcal{C}\!\mathit{omplexes}}
% \Pic is the operator that assigns to X its picard group, usage \Pic(X)
% \Picardstack_{X/B} denotes the Picard stack of X over B
% \Picardfunctor_{X/B} denotes the Picard functor of X over B
\def\Pic{\mathop{\mathrm{Pic}}\nolimits}
\def\Picardstack{\mathcal{P}\!\mathit{ic}}
\def\Picardfunctor{\mathrm{Pic}}
\def\Deformationcategory{\mathcal{D}\!\mathit{ef}}
\usepackage{cleveref}
\crefname{originalreference}{original reference}{original references}
\Crefname{originalreference}{Original reference}{Original references}
\makeatletter
\newlabel{lemma-cohomological-dimension-proper}{{095U}{source}{}{sourcepointer-1}{}}
\newlabel{lemma-cohomological-dimension-proper@cref}{{[originalreference][0][]095U}{[1][0][]source}{}{}{}}
\makeatother
\crefname{originalreference}{original reference}{original references}
\Crefname{originalreference}{Original reference}{Original references}
\makeatletter
\newlabel{lemma-proper-base-change-mod-n}{{0F0C}{source}{}{sourcepointer-2}{}}
\newlabel{lemma-proper-base-change-mod-n@cref}{{[originalreference][0][]0F0C}{[1][0][]source}{}{}{}}
\makeatother
\crefname{originalreference}{original reference}{original references}
\Crefname{originalreference}{Original reference}{Original references}
\makeatletter
\newlabel{lemma-pull-out-constant-mod-n}{{0F12}{source}{}{sourcepointer-3}{}}
\newlabel{lemma-pull-out-constant-mod-n@cref}{{[originalreference][0][]0F12}{[1][0][]source}{}{}{}}
\makeatother
\crefname{originalsection}{original section}{original sections}
\Crefname{originalsection}{Original section}{Original sections}
\makeatletter
\newlabel{sites-cohomology-section-projection-formula}{{0943}{source}{}{sourcepointer-4}{}}
\newlabel{sites-cohomology-section-projection-formula@cref}{{[originalsection][0][]0943}{[1][0][]source}{}{}{}}
\makeatother
\title{The projection formula for a proper morphism with coefficients in a torsion ring}
\author{The Stacks Project Authors}
\date{stacks-project:etale-cohomology.tex:local-snapshot-20261011; GNU FDL 1.2 or later}
\begin{document}
\maketitle
\noindent\textit{Editorial note.} All delivered mathematical text is the author's own, copied from the pinned Stacks Project chapter file etale-cohomology.tex; no statement or proof step was rewritten, reordered or removed.\par\medskip
\noindent\textit{Editorial note.} The counted proof is the author's complete argument: the proof is written out in full: the stalkwise check, the reduction by proper base change, and the application of the constant-factor lemma with the dimension bound.\par\medskip
\noindent\textit{Editorial note.} The label lemma-projection-formula-proper-mod-n names the counted statement, and the proof environment immediately following it in the source is delivered verbatim as its proof.\par\medskip
\noindent\textit{Editorial note.} The counted statement is carried by the author's own lemma environment in the chapter source.\par\medskip
\noindent\textit{Editorial note.} The statement and proof are copied from the pinned chapter file; the four cited inputs are kept as external pointers with their canonical Stacks tags.\par\medskip
\noindent\textit{Allowed standard prerequisites (editorial).}\begin{itemize}
\item The canonical projection map. Cohomology on Sites, the section on the projection formula providing the canonical map Rf\_*E tensor\textasciicircum{}L K -> Rf\_*(E tensor\textasciicircum{}L f\textasciicircum{}\{-1\}K). Supplies the map whose being an isomorphism is proved.
\item Proper base change for complexes of modules over a torsion ring. Same chapter, the proper base change theorem for modulo n coefficients. Supplies the reduction of the stalkwise check to the case of an algebraically closed field base.
\item Pulling the constant factor out for finite cohomological dimension. Same chapter, the lemma computing global sections of a tensor product with a constant complex, together with the finiteness of the cohomological dimension of a proper scheme. Supplies the case of an algebraically closed field base.
\end{itemize}\medskip
\ifdefined\hypertarget\hypertarget{sourcepointer-1}{}\fi
\noindent\textit{Source reference.} Original reference 095U: sources/stacks-etale-cohomology/pinned-original-source.tex, source lines 18786--18786. Prerequisite pointer invoked at proof line(s) 20214; it names a standard fact of the same chapter that the counted argument uses, and every step of the counted proof built on it is retained verbatim in the delivered body.\par\medskip
\ifdefined\hypertarget\hypertarget{sourcepointer-2}{}\fi
\noindent\textit{Source reference.} Original reference 0F0C: sources/stacks-etale-cohomology/pinned-original-source.tex, source lines 18887--18887. Prerequisite pointer invoked at proof line(s) 20209; it names a standard fact of the same chapter that the counted argument uses, and every step of the counted proof built on it is retained verbatim in the delivered body.\par\medskip
\ifdefined\hypertarget\hypertarget{sourcepointer-3}{}\fi
\noindent\textit{Source reference.} Original reference 0F12: sources/stacks-etale-cohomology/pinned-original-source.tex, source lines 20164--20164. Prerequisite pointer invoked at proof line(s) 20212; it names a standard fact of the same chapter that the counted argument uses, and every step of the counted proof built on it is retained verbatim in the delivered body.\par\medskip
\ifdefined\hypertarget\hypertarget{sourcepointer-4}{}\fi
\noindent\textit{Source reference.} Original section 0943: sources/stacks-chow/tags.txt, source lines 11536--11536. Prerequisite pointer invoked at proof line(s) 20206; it names a standard fact of a sibling chapter that the counted argument uses, and every step of the counted proof built on it is retained verbatim in the delivered body.\par\medskip
\begin{lemma}
\label{lemma-projection-formula-proper-mod-n}
Let $f : X \to Y$ be a proper morphism of schemes. Let $\Lambda$
be a torsion ring. Let $E \in D(X_\etale, \Lambda)$ and
$K \in D(Y_\etale, \Lambda)$. Then
$$
Rf_*E \otimes_\Lambda^\mathbf{L} K =
Rf_*(E \otimes_\Lambda^\mathbf{L} f^{-1}K)
$$
in $D(Y_\etale, \Lambda)$.
\end{lemma}

\begin{proof}
There is a canonical map from left to right by
Cohomology on Sites, Section \ref{sites-cohomology-section-projection-formula}.
We will check the equality on stalks at $\overline{y}$.
By the proper base change (in the form of
Lemma \ref{lemma-proper-base-change-mod-n} where $Y' = \overline{y}$)
this reduces to the case where $Y$ is the spectrum of
an algebraically closed field.
This is shown in Lemma \ref{lemma-pull-out-constant-mod-n}
where we use that $\text{cd}(X) < \infty$ by
Lemma \ref{lemma-cohomological-dimension-proper}.
\end{proof}

\end{document}
